Una espira rectangular i conductora es troba a prop d'un fil conductor rectilini infinit pel qual circula una intensitat de corrent $I$ cap avall, tal com es mostra en la figura \textit{a}.

\begin{center}
\begin{minipage}[b]{0.3\linewidth}
\figura[1]{fig1.png}
{\centering (\textit{a})\par}
\end{minipage}\hspace{0.06\linewidth}%
\begin{minipage}[b]{0.48\linewidth}
\figura[1]{fig2.png}
{\centering (\textit{b})\par}
\end{minipage}
\end{center}

\begin{apartat}{1}
Representeu el sentit i la direcció del camp magnètic creat pel fil conductor en la regió plana delimitada per l'espira. Aquest camp magnètic és uniforme en la regió delimitada per l'espira? Justifiqueu la resposta.
\end{apartat}

\begin{apartat}{1}
El fil conductor i l'espira no es mouen, però la intensitat del corrent que circula pel conductor varia amb el temps, tal com indica la gràfica (figura \textit{b}). Argumenteu si s'indueix o no corrent en l'espira en els intervals de temps següents: de 0 a 20~s, de 20 a 80~s i de 80 a 120~s. En quin d'aquests tres intervals de temps la intensitat del corrent induït és més gran? Justifiqueu la resposta.
\end{apartat}
